\section*{Chapter \ref{ch:b2:charges-currents-conduction} --- Charges, Currents and Conduction}

\begin{solution}{exo:b2:charges-currents-conduction:1}
$j = 16/2.5 \times 10^{-6} = \qty{6.4e6}{A/m^2}$; $v = j/ne = \qty{0.47}{mm/s}$; \qty{20}{m}
in $4.3 \times 10^4$ s, twelve hours. $E = j/\gamma = \qty{0.11}{V/m}$; $U = \qty{2.1}{V}$;
$P = UI = \qty{34}{W}$.
\end{solution}

\begin{solution}{exo:b2:charges-currents-conduction:2}
Lightning $3 \times 10^4/\pi10^{-4} = \qty{1e8}{A/m^2}$; beam $\qty{4e3}{A/m^2}$; axon
$10^{-9}/1.3 \times 10^{-11} = \qty{80}{A/m^2}$; ring current \qty{1e-8}{A/m^2}. Only the
lightning exceeds the wire's $6 \times 10^6$.
\end{solution}

\begin{solution}{exo:b2:charges-currents-conduction:3}
$R = \rho_eL/S$: $0.68$, $1.1$, $4.0$, \qty{44}{\ohm}. Aluminium: $2.5 \times 2.8/1.7 =
\qty{4.1}{mm^2}$, mass ratio $(4.1 \times 2700)/(2.5 \times 8900) = 0.50$ --- half the mass.
Heater: $R = 230^2/1000 = \qty{53}{\ohm}$, $S = \qty{0.196}{mm^2}$, $L = RS/\rho_e = \qty{9.4}{m}$.
\end{solution}

\begin{solution}{exo:b2:charges-currents-conduction:4}
$\tau = m\gamma/ne^2$: copper \qty{2.5e-14}{s}, $\mu = e\tau/m = \qty{4.4e-3}{m^2/(V.s)}$,
$\ell = \qty{40}{nm}$; silver \qty{3.8e-14}{s}, \qty{6.7e-3}{m^2/(V.s)}, \qty{61}{nm}.
Silicon: $\gamma = ne\mu = 10^{22} \times 1.6 \times 10^{-19} \times 0.14 = \qty{220}{S/m}$; $v =
\mu E = \qty{140}{m/s}$.
\end{solution}

\begin{solution}{exo:b2:charges-currents-conduction:5}
(a) Current in, none out: $\dd Q/\dd t = I$. (b) Charge accumulates: $\partial_t
\rho \ne 0$. (c) $\vect j = (I/4\pi r^2)\vect e_r$: $\operatorname{div}\vect j = \frac1{r^2}
\partial_r(I/4\pi) = 0$. (d) Only at the electrode, where the current is
injected from the wire --- a source of the flux.
\end{solution}

\begin{solution}{exo:b2:charges-currents-conduction:6}
(a) $j = I/2\pi rL$, $E = j/\gamma$, radial. (b) $U = \int_a^bE\,\dd r = I\ln(b/a)/
2\pi\gamma L$: $R = \ln3.5/(2\pi \times 10^{-13} \times 10^3) = \qty{2e9}{\ohm}$ per kilometre.
(c) \qty{0.5}{\micro A}, \qty{0.5}{mW}. (d) $RC = \varepsilon_0\varepsilon_r/\gamma = \qty{200}{s}$ for
$\varepsilon_r = 2.3$: the cable, charged and isolated, discharges through its own
insulator with the relaxation time of that insulator.
\end{solution}

\begin{solution}{exo:b2:charges-currents-conduction:7}
(a) $\omega\tau$: $8 \times 10^{-12}$, $1.6 \times 10^{-3}$, $4.7$, $79$. (b) $|\underline\gamma|/\gamma_0 =
1/\sqrt{1 + \omega^2\tau^2}$: $1$, $1$, $0.21$, $0.013$; phase $-\arctan\omega\tau$: $0$,
$-\qty{0.09}{\degree}$, $-\qty{78}{\degree}$, $-\qty{89}{\degree}$. (c) $\omega\tau \le 0.14$: $f \le
\qty{0.9}{THz}$. (d) $\underline\gamma \approx ne^2/\iu m\omega$: $\underline{\vect j}$ lags
$\underline{\vect E}$ by \qty{90}{\degree}, $\langle\vect j\cdot\vect E\rangle = 0$: no Joule
heating --- the electrons oscillate freely, as in a plasma.
\end{solution}

\begin{solution}{exo:b2:charges-currents-conduction:8}
(a) $S = \qty{3.1e-8}{m^2}$, $j = \qty{3.2e8}{A/m^2}$, $p = j^2/\gamma = \qty{1.7e9}{W/m^3}$.
(b) Energy to melt per unit volume $\rho(c\Delta T + L_f) = 8900(4.1 \times 10^5 + 2.05 \times
10^5) = \qty{5.5e9}{J/m^3}$: \qty{3.2}{s}. (c) Nine times faster: \qty{0.36}{s}. (d) The
rated current is the one whose Joule heating the air just evacuates
below the melting point; the sand quenches the arc that would
otherwise keep conducting across the gap after the wire melts.
\end{solution}

\begin{solution}{exo:b2:charges-currents-conduction:9}
(a) $\partial_t\rho = -\operatorname{div}\vect j = -\gamma\operatorname{div}\vect E = -\gamma\rho/
\varepsilon_0$. (b) $\varepsilon_0/\gamma$: \qty{1.5e-19}{s}, \qty{1.8e-12}{s}, \qty{1.8e-10}{s},
\qty{9}{s}, $9 \times 10^4$ s. (c) $f \lesssim 1/2\pi\tau_r$: $10^{18}$ Hz (the \omterm{prop:b2:charges-currents-conduction:drude}{Drude model}
fails long before), \qty{90}{GHz}, \qty{0.9}{GHz}, \qty{0.02}{Hz}, \qty{2e-6}{Hz}. (d)
A glass rod would lose its charge in a minute through its volume;
real dry glass is nearer $10^{-14}$ S/m (hours), and it is surface
humidity that usually discharges it --- consistent in order of
magnitude.
\end{solution}

\begin{solution}{exo:b2:charges-currents-conduction:10}
(a) $\vect j = (I/2\pi r^2)\vect e_r$, $E = I/2\pi\gamma r^2$, $V = I/2\pi\gamma r$. (b) $R = V(a)
/I = 1/2\pi\gamma a = \qty{32}{\ohm}$. (c) $V_{\text{rod}} = \qty{640}{kV}$; $V(10) = \qty{32}{kV}$,
$V(10.7) = \qty{30}{kV}$: \qty{2.1}{kV} between the feet. (d) Front and hind
legs \qty{1.5}{m} apart along the radius: $4$--\qty{5}{kV}, and the current
path crosses the heart.
\end{solution}

\begin{solution}{exo:b2:charges-currents-conduction:11}
(a) $E_H = jB/nq$ with $j = I/wt$: $V_H = E_Hw = IB/nqt$. (b) $1/(8.5 \times 10^{28}
\times 1.6 \times 10^{-19} \times 10^{-4}) = \qty{0.7}{\micro V}$. (c) \qty{6.3}{mV}: $V_H \propto 1/n$,
and semiconductors have $10^6$ times fewer carriers. (d) $B = V_Hnqt/I =
10^{-5} \times 10^{22} \times 1.6 \times 10^{-19} \times 2 \times 10^{-5}/10^{-3} = \qty{0.32}{T}$; the sign
of $V_H$ gives the sign of the carriers (for a known $\vect B$), or the
direction of $\vect B$ (for a known sensor).
\end{solution}

\begin{solution}{exo:b2:charges-currents-conduction:12}
(a) $\vect j = ne\mu_+\vect E + (-ne)(-\mu_-\vect E)$: both species drift in
opposite directions and carry current in the same direction. (b) $n =
0.5 \times 10^3 \times 6.02 \times 10^{23} = \qty{3e26}{m^{-3}}$: $\gamma = 3 \times 10^{26} \times 1.6 \times 10^{-19}
\times 1.31 \times 10^{-7} = \qty{6.3}{S/m}$ (the ions hinder each other at this
concentration: \qty{5}{S/m} measured). (c) \qty{0.52}{\micro m/s} and \qty{0.79}{\micro m/s}:
Cl$^-$ carries $60\%$. (d) $R = L/\gamma S = \qty{200}{\ohm}$, $U = \qty{200}{V}$, $P =
\qty{200}{W}$ into \qty{10}{g} of water: \qty{10}{K} in \qty{2}{s} --- hence the cooling
plates of electrophoresis tanks.
\end{solution}

\begin{solution}{pb:b2:charges-currents-conduction:1}
\textbf{1.} $n = \rho N_A/M = \qty{8.4e28}{m^{-3}}$.

\textbf{2.} $j = \qty{6.4e6}{A/m^2}$, $v = j/ne = \qty{0.47}{mm/s}$; \qty{35}{min} per metre.

\textbf{3.} $\tau = m\gamma/ne^2 = \qty{2.5e-14}{s}$, $\mu = \qty{4.4e-3}{m^2/(V.s)}$, $\ell =
\qty{40}{nm}$, $150$ spacings: the electrons collide with defects and
thermal vibrations, not with the atoms.

\textbf{4.} $\rho_e \propto T$ gives $\alpha = 1/293 = \qty{3.4e-3}{K^{-1}}$, close to the
measured $4 \times 10^{-3}$.

\textbf{5.} $E = \qty{0.11}{V/m}$, $U = \qty{0.11}{V}$; $\tfrac12mv^2 = \qty{1e-37}{J}$ against
$k_BT = \qty{4e-21}{J}$: the electron gas does not notice the current.

\textbf{6.} $I/e = \qty{1e20}{}$ electrons per second.

\textbf{7.} $Q = neSL = \qty{3.4e4}{C}$; $QE = \qty{3.6}{kN}$, against a weight of
\qty{0.2}{N}. The field pulls the positive lattice with the opposite
force; the electrons hand their momentum to the lattice by collisions:
the wire as a whole feels nothing.

\textbf{8.} $p = j^2/\gamma = \qty{6.8e5}{W/m^3}$, \qty{1.7}{W/m}.

\textbf{9.} Surface $\pi d = \qty{5.6e-3}{m^2}$ per metre: $\Delta T = 1.7/(10 \times 5.6 \times
10^{-3}) = \qty{30}{K}$; the resistivity rises $12\%$, the power to \qty{1.9}{W}:
$\Delta T \approx \qty{35}{K}$.

\textbf{10.} $\Delta T = 1.7(1 + 0.004\Delta T)/(2 \times 5.6 \times 10^{-3})$: $\Delta T(1 - 0.6) = 150$,
$\Delta T \approx \qty{400}{K}$ --- near thermal runaway: buried cables must be
derated.

\textbf{11.} $p/\rho c = \qty{0.2}{K/s}$; at \qty{160}{A}, \qty{20}{K/s}: melting in
about a minute (less as $\rho_e$ rises); the breaker opens in milliseconds.

\textbf{12.} $\omega\tau = 8 \times 10^{-12}$: Ohm's law is untouched.

\textbf{13.} $\varepsilon_0/\gamma = \qty{1.5e-19}{s}$: no charge inside; the field in the
wire is made by charges on its \emph{surface}, whose density varies
along the wire.

\textbf{14.} $Q = I_0\tau_c = \qty{16}{mC}$, $V = Q/C = \qty{160}{V}$.

\textbf{15.} The wire side carries $I$, the gap side no conduction
current: the charge grows, and with it the electric field in the gap.

\textbf{16.} $I = 16\eu^{-0.5} = \qty{9.7}{A}$; $E = Q/\varepsilon_0A$, $\dd E/\dd t = I/\varepsilon_0A
= 9.7/(8.85 \times 10^{-12} \times 10^{-4}) = \qty{1.1e16}{V.m^{-1}.s^{-1}}$.

\textbf{17.} $\varepsilon_0A\,\dd E/\dd t = I$: the displacement current.

\textbf{18.} $j = \gamma E = 10^{-14} \times 2.3 \times 10^5 = \qty{2e-9}{A/m^2}$, $10^{-14}$ A
through the section: nothing; breakdown at \qty{3}{kV} across the millimetre.

\textbf{19.} $R = 1000/(6 \times 10^7 \times 10^{-5}) = \qty{1.7}{\ohm}$; $\Delta U = \qty{67}{V}$; $P =
\qty{2.7}{kW}$, $29\%$ of the \qty{9.2}{kW} delivered.

\textbf{20.} $I = \qty{0.46}{A}$, $\Delta U = \qty{0.8}{V}$, \qty{0.35}{W}: transmission at
high voltage divides the loss by $(U_2/U_1)^2$.

\textbf{21.} $10 \times 6/3.6 = \qty{17}{mm^2}$; \qty{89}{kg/km} of copper against
\qty{45}{kg/km} of aluminium: lighter and cheaper on pylons; in the wall
copper's smaller section fits the conduits and terminals.

\textbf{22.} $V_H = IB/nqt = 5 \times 10^{-3} \times 0.1/(10^{22} \times 1.6 \times 10^{-19} \times 10^{-4}) =
\qty{3.1}{mV}$.

\textbf{23.} $+24\%$: \qty{2.1}{\ohm}, \qty{3.3}{kW} lost.

\textbf{24.} $RI^2t = 1.7 \times 4 \times 10^8 \times 10^{-4} = \qty{67}{kJ}$ into \qty{89}{kg}: \qty{2}{K}.
$j = \qty{2e9}{A/m^2}$, three hundred times the household value.

\textbf{25.} $n$ (matter) and $\tau$ (collisions) give $\gamma = ne^2\tau/m$;
$\gamma$ links $j$ and $E$ (Ohm); $j$ integrates to $I$, $E$ to $U$ (the
circuit); $j^2/\gamma$ is the heat.
\end{solution}
